\section{Signed Tait block products for actual observer workloads}
\label{sec:tait-blocks}

The preceding boundary experiment isolates a useful primitive, but actual
recognition decides when that primitive runs.  We therefore traced the
unmodified marked observer on every retained example before choosing its
next optimization.  Most examples ask for just the initial determinant.
The determinant-one torus and grid examples ask for two additional distinct
closures at a later stage.  No example in this collection reaches the
repeated-query threshold of four at any stage.  This evidence favors reducing
the cost of a single completed Tait graph before enabling terminal-kernel
setup in the production observer.

The maintained shadow backend now factors larger signed Tait graphs at their
articulation vertices.  For example, the initial eightfold Conway connected
sum has a $41\times41$ full cofactor, but its block factors require matrices
of size at most $6\times6$.  This optimization applies directly to each
actual completion, with the original signed phase retained.  It makes no
assumption that a partial complex splits as a knot connected sum.

\subsection{Why the signed product identity is exact}

Let $G$ be the black Tait multigraph of a connected classical projection,
with integer edge weights $w_e\in\{1,-1\}$.  Its weighted tree sum is
\[
 \tau(G)=\sum_{T\text{ spanning tree of }G}\prod_{e\in T}w_e.
\]
Loops never occur in a tree and may be discarded for this computation.
For parallel edges between $u$ and $v$, a tree can use at most one, so
replacing them by the single weight
$w_{uv}=\sum_{e:u\leftrightarrow v}w_e$ preserves the sum.  This is a
polynomial identity and does not require positive weights.  A zero combined
weight may be deleted.  If these deletions disconnect the graph, the sum is
zero even though the original knot projection remains connected.

Suppose the resulting graph is connected, and let $G_1,\ldots,G_s$ be its
edge blocks, with each bridge treated as a two-vertex block.  Then
\[
 \boxed{\tau(G)=\prod_{j=1}^s\tau(G_j).}
\]
To prove this, first split at an articulation vertex into two graphs meeting
only at that vertex.  Restricting a spanning tree to either graph gives a
spanning tree there: a disconnected restriction cannot be reconnected by
leaving and returning through the same articulation vertex.  Conversely,
the union of two such trees is connected and has no cycle.  The bijection
preserves the product of edge weights.  Induction on the block-cut tree
gives the formula.  Thus the identity also holds when factors are negative
or vanish by cancellation.  A bridge contributes its combined weight; the
one-vertex graph has tree sum one.

For a nontrivial block, the integer matrix-tree theorem identifies its tree
sum with any grounded Laplacian cofactor.  The implementation uses the
existing fraction-free Bareiss determinant with exact-division checks and
row-level interruption points.  It stops once a zero factor is found.
No positive-definite elimination or nonsingular-interior premise is used.

The full specialization is still
\[
 J_{\mathrm{raw}}(i)=(-i)^{B+v-1}\prod_j\tau(G_j),
\]
where $v$ is the black vertex count of the \emph{full completed graph} and
$B$ is the original crossing sum.  In particular, loops and cancelled
parallel edges continue to contribute to the crossing phase.  Multiplying
independently normalized block Jones polynomials without these conventions
would not justify the implementation.

This factorization is performed after the actual closure has been checked.
It cannot automatically be applied to the fixed common graph before arbitrary
terminal identifications: those identifications can create cycles crossing
old articulation blocks.  Terminal-kernel reuse and completed-graph block
factorization therefore have different geometric premises.

\subsection{Implementation, resources, and complexity}

\path{fast/fastunknot/tait_blocks.py} combines parallel weights and uses an
iterative low-link depth-first search to produce edge blocks.  A stack of
DFS frames replaces recursion, so a long chain does not impose Python's
recursion-depth limit.  When a child's low-link index is at least its
parent's discovery index, the edge stack is popped through that tree edge
to emit a block.  A visit count detects disconnection of the nonzero graph.

The observer chooses the smaller checkerboard shading as before.  For a
cofactor of size at most sixteen it retains the direct dense path; larger
cofactors use block decomposition before any dense matrix is allocated.
This cutoff avoids adding graph bookkeeping to the smallest queries.  It
is an implementation policy, not a claim that sixteen is universally optimal.
The observer exposes block, bridge, cancellation, and block-determinant
counts.  \code{max\_unsplit\_cofactor\_size} records the full graph's size,
while \code{max\_cofactor\_size} now records the largest matrix actually
allocated.  The preexisting \code{determinants} field continues to count
connected closure specializations, including interrupted attempts.

Write $b_j=|V(G_j)|-1$.  Apart from integer bit costs, the determinant path
uses
\[
 O\!\left(|V|+|E|+\sum_j b_j^3\right)
 \quad\text{work and}\quad
 O\!\left(|V|+|E|+\max_j b_j^2\right)\text{ auxiliary space},
\]
with bridge factors evaluated directly.  This replaces the full cofactor's
cubic arithmetic and quadratic allocation by the corresponding block costs.
The graph decomposition itself is linear.  Integer sizes still grow with
the input; the displayed arithmetic count is not a unit-cost claim for
arbitrarily large numbers.  These polynomial savings do not bound the
Khovanov scanner's exponential state growth or establish a general
subexponential recognition algorithm.

All graph traversal, edge aggregation, matrix allocation, and elimination
paths charge the existing local work allowance and poll the global deadline.
Only a complete four-residue result enters the closure cache.  A partially
multiplied block product is local state and is discarded on interruption.
Local exhaustion disables optional observations and continues the same
exact capped scan; a global deadline still yields \code{UNKNOWN}.  The
bounded-work policy counts combinatorial and arithmetic updates, not integer
bit operations.  The change remains inside the optional \code{shadow}
backend; default recognition still uses \code{standard}.

\subsection{Validation and workload evidence}

Five new tests compare 245 random small weighted multigraphs with independent
spanning-tree enumeration and 192 larger graphs with independent dense
cofactors.  They include parallel edges, loops, negative weights, deletion
of cancelled edges, and connected graphs whose signed tree sum vanishes.
A 5,000-vertex bridge chain verifies the iterative traversal and avoids
any dense allocation.  Four larger classical diagrams and their mirrors
match independently reconstructed values at both one and $i$; this checks
that the block product retains the global phase.

An injected interruption after one completed block verifies that no partial
result is cached.  Resumption recomputes the exact value, global deadlines
win even on a cache hit, and a real 2,000-unit optional allowance on the
Conway sum exhausts locally while exact scanning still returns knotted.
All 577 maintained tests pass after the final cutoff change.
The test records are \path{data/tait-block-tests.txt} and
\path{data/tait-block-integrated-tests.txt}.

\path{fast/benchmark_tait_blocks.py} keeps three scopes separate: the initial
closure query, the complete raw marked scan, and complete recognition with
the shadow backend requested.  Each sample includes fresh diagram and
observer setup.  Two additional inputs concatenate 16 and 32 oriented Conway traversals to
form larger connected sums; their validated PD codes are retained in the
benchmark data.  Seven paired rounds shuffle direct, identical control, and
block-factor arms; imports and one warmup are excluded.  The direct baseline
loads \code{shadow\_scan.py} from revision \code{2c7c8df3f}, sharing other
dependencies.  Actual observer traces are collected separately outside the
timed calls.  Inputs, source hashes, work counts, and raw samples accompany
the resulting tables.

\begin{center}\small
\begin{tabular}{@{}lrrrrrr@{}}
\toprule
Input & Initial & A/A & Raw & A/A & Full & A/A \\
\midrule
Conway & 0.983 & 0.994 & 1.004 & 0.999 & 0.999 & 0.999 \\
Conway double three & 0.999 & 0.996 & 1.007 & 1.005 & 0.998 & 0.998 \\
Conway pretzel & 1.043 & 1.061 & 0.972 & 0.995 & 0.967 & 0.992 \\
Conway sum 2 & 0.995 & 1.000 & 1.015 & 0.998 & 0.995 & 0.993 \\
Conway sum 3 & 1.008 & 1.009 & 0.984 & 0.987 & 1.048 & 1.001 \\
Conway sum 8 & 2.152 & 0.999 & 1.186 & 1.009 & 1.030 & 1.027 \\
Figure eight & 1.001 & 1.026 & 1.009 & 1.013 & 1.003 & 0.994 \\
Grid determinant one & 0.991 & 0.982 & 0.999 & 0.993 & 1.005 & 0.989 \\
Scrambled unknot & 0.984 & 1.006 & 0.986 & 0.998 & 0.996 & 0.997 \\
Hard unknot 8 & 0.994 & 0.997 & 1.018 & 1.013 & 1.003 & 0.994 \\
Kinoshita--Terasaka & 0.996 & 1.004 & 0.987 & 0.991 & 1.001 & 1.004 \\
Stress braid 36 & 1.004 & 1.000 & 0.998 & 0.986 & 0.986 & 0.992 \\
$T(3,5)$ & 0.994 & 1.007 & 0.999 & 0.997 & 0.999 & 0.990 \\
Trefoil & 0.989 & 0.996 & 0.982 & 0.995 & 0.992 & 1.007 \\
Crossingless unknot & --- & --- & 0.986 & 1.002 & 0.983 & 1.003 \\
Braid unknot 40 & 1.302 & 0.989 & 1.082 & 1.009 & 1.000 & 1.002 \\
Conway sum 16 & 5.236 & 1.007 & 1.623 & 1.008 & 1.000 & 0.977 \\
Conway sum 32 & 17.418 & 0.995 & 21.547 & 1.024 & 1.001 & 1.000 \\
\bottomrule
\end{tabular}
\end{center}
Ratios are median paired direct/block times; values above one favor block evaluation. A/A compares identical direct implementations. Every scope includes fresh setup.


\paragraph{Choosing the direct-path cutoff from measured overhead.}
The first run, preserved in
\path{fast/results/tait_blocks_threshold8_20261008.json}, used cutoff eight
and the implementation at \code{07c02e10c}.  It improved the eightfold Conway
sum's initial query by $2.10$ times, but made the irreducible size-16 stress
cofactor about nine percent slower.  The size-11 twofold sum also lost about
five percent.  Raising the cutoff to sixteen removes those decomposition
costs, at the price of foregoing the initial run's modest size-16 threefold
sum gain.  The displayed final run uses this revised policy.  Both source
versions and all measurements are retained; the comparison does not replace
the earlier unfavorable samples with the later ones.

\paragraph{Single queries and complete raw scans.}
The final initial-query speedups are $2.15$, $5.24$, and $17.42$ times for
8, 16, and 32 Conway summands.  The 40-crossing braid unknot improves by
$1.30$ times: all twenty effective blocks are bridges, so it allocates no
dense determinant matrix.  The size-16 irreducible stress case now keeps
its direct path and has ratios near one.  The corresponding complete raw
scan gains are $1.19$, $1.62$, and $21.55$ times for the three larger sums,
and $1.08$ times for the braid unknot.

The 32-fold sum illustrates why resource behavior matters.  Its initial
unrestricted direct query needs 1,398,785 work units, compared with 13,252
for block products.  With the normal million-unit optional allowance,
the old driver stops observing during this query and scans all 352
crossings, eventually certifying knottedness from closed rank.  The new
driver uses only 14,619 units including its later Euler query and obtains
a component-Euler certificate after 11 crossings.  Its largest determinant
is $5\times5$, compared with the original $160\times160$ cofactor.
The raw benchmark therefore requires agreement of verdicts while recording
method and stage differences explicitly.  The large gain here includes
preserving the opportunity for later inference, not just faster arithmetic.
Both runs complete; no timeout is assigned a speedup ratio.

\paragraph{Complete recognition and the remaining limit.}
Earlier recognizer filters decide all 18 measured inputs before the shadow
backend runs.  Full-recognition ratios range from $0.967$ to $1.048$, with
no demonstrated overall benefit from the new determinant path.  For the
32-fold sum the existing factorization and Jones filter already return in
about 15 milliseconds, ahead of either raw scanner.  The results support
the optimization inside the optional observer; they do not justify changing
the default backend or claiming a general asymptotic improvement in unknot
recognition.  Final inputs, traces, source hashes, and measurements are in
\path{fast/results/tait_blocks_20261008.json}.
